\documentclass{article}[17pt]

\usepackage{../uni-notes-eng}
\usepackage{multicol}
\usepackage{graphicx}
\usepackage{tabularx}
\usepackage{colortbl}

\arrayrulecolor{sep}

\title {Kombinatoryka $\&$ teoria grafów}
\author{by a fish}
\date {21.03.2137}

\begin{document}

\maketitle

\newpage

\begin{center}
    \includegraphics[height=\textheight]{potega.jpg}
\end{center}

\newpage

\section*{SYLABUS - teoria grafów:}
1. Basic concepts: graphs, paths and cycles, complete andbipartite graphs\\
2. Matchings: Hall's Marriage theorem and its variations\\
3. Forbidden subgraphs: complete bipartite and r-partite subgraphs, chromatic numbers, Tur"an's thorem, asymptotic behaviour og edge density, Erd"os-Stone theorem\\
4. Hamiltonian cycles (Dirac's Theorem), Eulerian circuits\\
5. Connectivity: connected and k-connected graphs, Menger's theorem\\
6. Ramsey theory: edge colourings of graphs, Ramsey's theorem and its variations, asymptotic bounds on Ramsey numbers\\
7. Planar graphs and colourings: statements of Kuratowski's and Four Colour theorems, proof of Five Colour theorem, graphs on other surfaces and Euler chracteristics, chromatic polynomial, edge colourings and Vizing's theorem\\
8. Random graphs: further asymptotic bounds on Ramsey numbers, Zarankiewicz numbers and their bounds, graphs of large firth and high chromatic number, cmplete subgraphs in random graphs.\\
9. Algebraic methods: adjavenvy matrix and its eigenvalues, strongly regular graphs, Moore graphs and their existence.

\newpage

\tableofcontents

\newpage

\section{Teoria grup}

\subsection{Grupy, pierscienie, ciala}
\begin{multicols}{2}

    Dzialanie na zbiorze $X$:
    $$\Phi: X\times X\to X,$$
    zwykle zapisywane jako $xy$, $x\cdot y$, $x+y$.\smallskip

    {\color{def}Element neutralny} - takie $e$, ze dla kazdego $x\in X$ $ex=xe=x$. Dzialanie ma co najwyzej jeden element neutralny.\smallskip
    
    {\color{def}Element odwrotny} do $x$ to takie $y$, ze $xy=yx=e$. Jesli dzialanie jest laczne, to ma co najwyzej jeden element odwrotny do danego $x$.\medskip

    \podz{sep}\medskip

    {\color{def}Homomorfizm} algebry $\rodz X=(X,\{ \cdot\})$ na algebre $\rodz Y=(Y,\{\circ\})$ nazywamy przeksztalcenie $f:X\to Y$ spelniajace dla kazdego $a, b\in X$
    $$f(a\cdot b)=f(a)\circ f(b).$$
    \indent {\color{acc}$\bullet$ monomorfizm} - $f$ jest 1-1\smallskip\\
    \indent {\color{acc}$\bullet$ epimorfizm} - $f$ jest "na"\smallskip\\
    \indent {\color{acc}$\bullet$ izomorfizm} - $f$ jest 1-1 i "na"\smallskip\\
    \indent {\color{acc}$\bullet$ endomorfizm} - kiedy $\rodz Y=\rodz X$\smallskip\\
    \indent {\color{acc}$\bullet$ automorfizm} - enodmorfizm bedacy izomorfizmem\medskip

    \podz{sep}\medskip

    {\color{def}Polgrupa} to niepusty zbior z dzialaniem lacznym.\medskip
    
    {\color{def}GRUPA} to niepusty zbior z lacznym dzialaniem i elementem neutralnym (zwanym {\color{acc}jednoscia grupy}) oraz elementami odwrotnymi dla kazdego elementu.\smallskip\\
    \indent {\color{acc}$\hookrightarrow$ grupa abelowa} (przemienna) - grupa z dzialaniem przemiennym\smallskip\\
    Zbior $G$ z dzialaniem $\cdot$ jest grupa, jesli:\smallskip\\
    \indent 1. $(\forall\;a,b,c\in G)\;(ab)c=a(bc)$\\
    \indent 2. $(\exists\;e\in G)(\forall\;a\in G)\;ea=ae=e$\\
    \indent 3. $(\forall\;a\in G)(\exists\;b\in G)\;ab=ba=e$\\
    \indent *4. $(\forall\;a,b\in G)\;ab=ba$ w grupie \emph{abelowej}\bigskip

    {\color{def}PIERSCIEN} to niepusty zbior $X$ z dwoma dzialaniami ($\cdot,\;+$, mnozenie i dodawanie) taki, ze:\smallskip\\
    \indent 1. zbior $X$ z $+$ jest grupa abelowa\smallskip\\
    \indent 2. $\cdot$ jest laczne\smallskip\\
    \indent 3. $(\forall\;x,y,z\in X)\;x\cdot(y+z)=x\cdot y+x\cdot z\;\land\;(x+y)\cdot z=x\cdot z+y\cdot z$\smallskip\\
    Kolejne dzikie nazwy $\star$:\smallskip\\
    \indent {\color{acc}$\star$ pierscien przemienny} - jesli mnozenia jest przemienne\smallskip\\
    \indent {\color{acc}$\star$ pierscien z jednoscia} - dla mnozenia istnieje element neutralny\medskip

    {\color{def}CIALO} to pierscien przemienny, ktory dla kazdego elementu $\neq 0$ ma element odwrotny

\end{multicols}\bigskip

\podz{sep2}\bigskip

\subsection{Wlasnosci}
\begin{multicols}{2}

    Niech $G$ bedzie grupa, a $e$ jej elementem neutralnym. Wowczas:\smallskip\\
    \point $a,b\in G\implies (ab)^{-1}=b^{-1}a^{-1}$\smallskip\\
    \point $a\in G$ i $n=1, ..., n$ $a^{-n} = (a^n)^{-1} =^* (a^{-1})^n$\smallskip\\
    \point dla $m,n\in\Z$ i $a\in G$ mamy $a^{mn}=^* (a^m)^n$\smallskip\\
    \point dla $G$ grupy abelowej i $n\in\Z$ $(ab)^n=^*a^nb^n$\smallskip\\
    $^*$ trzeba udowodnic, ale mi sie nie chce\medskip

    $H\subseteq G$ jest {\color{def}podgrupa} $G$, jesli jest grupa ze wzgledu na te same dzialania, czyli wystarczy, ze 
    $$(\forall\;a,b\in H)\; ab^{-1}\in H.$$

    \podz{sep}\medskip

    Jelsi $a\in G$ i istnieja $n\in\N$, $n\geq 1$, takie, ze $a^n=e$, to mowimy ze $n$ jest {\color{def}rzedem elementu} $a$ ($n=o(a)$). Jesli takie $n$ nie istnieja, to $a$ ma {\color{acc}rzad nieskonczony} ($o(a)=\infty$).\smallskip\\
    {\color{acc}\point} {\color{def}grupa torsyjna} - wszystkie elementy maja rzad skonczony\smallskip\\
    {\color{acc}\point} {\color{def}grupa beztorsyjna} - wszystkie elementy maja rzad nieskonczony\smallskip\\
    \emph{Jesli $o(a) = n$ oraz $a^N=e$ to $n|N$, fajny dowodzik, ale leniem jestem}\medskip

    {\color{def}Grupa cykliczna} to grupa zlozona z wszystkich poteg danego elementu $a$, natomiast $a$ jest nazywane {\color{acc}generatorem} tej grupy

\end{multicols}
\newpage

\subsection{Permutacje :>}
\begin{multicols}{2}
    
    {\color{def}$n$-ta grupa symetryczna} [$S_n$] - grupa wszystkich permutacji zbioru $X_n = \{1, ..., n\}$. $|S_n|=n!$\medskip

    Jesli $P\in S_n$ i dla $i=1,...,n$ $P(i)=a_i$, to piszemy
    $$
        \begin{pmatrix}
            1   && 2   && ... && n  \\
            a_1 && a_2 && ... && a_n
        \end{pmatrix}
    $$

    Mnozenie permutacji:
    $$
        \begin{pmatrix}
            a_1 && a_2 && ... && a_n \\
            b_1 && b_2 && ... && b_n
        \end{pmatrix}
        \begin{pmatrix}
            c_1 && c_2 && ... && c_n \\
            a_1 && a_2 && ... && a_n
        \end{pmatrix} =
    $$
    $$=
    \begin{pmatrix}
        c_1 && c_2 && ... && c_n \\
        b_1 && b_2 && ... && b_n
    \end{pmatrix}
    $$

    {\color{def}Zbior elementow niezmienniczych} (fixpunktow) permutacji $P$ to zbior $F(P)=\{k\in X_n\;:\;P(k)=k\}$. Jego dopelnienie oznaczamy $M(P)=S_n\setminus F(P)$.\medskip

    \podz{sep}\medskip

    {\color{def}Cykl $k$-elementowy} $C$ to permutacja taka, ze $C(a_1)=a_2$, $C(a_2)=a_3$, ..., $C(a_n)=a_1$. Cykl $2$-elementowy to {\color{acc}transpozycja}. Cykle zapisujemy
    $$\begin{pmatrix}a_1, a_2,..., a_n\end{pmatrix}$$

    Kazda permutacja jest iloczynem transpozycji.\medskip

    \podz{sep}\medskip

    {\color{def}Permutacje parzyste} - iloczyn 
    $$\prod _{i<j}(a_j-a_i)$$ 
    jest dodatni (gorny row to kolejne liczby naturalne, dolny to wyrazy). Pozostale permutacje sa {\color{def}nieparzyste}.\medskip

    {\color{def}Znak permutacji} jest $+1$ gdzy permutacja jest parzysta i $-1$ wpp. Alternatywnie mozna zapisac (gorny row to $b_k$, a dolny to $c_k$)
    $$sgn\;P=\prod_{i<j}{b_j-b_i\over c_j-c_i}$$

    Dla dwoch dowolnych permutacji $P_1,P_2$ mamy
    $$sgn\;P_1P_2=sgn\;P_1\cdot sgn\;P_2$$
    $$sgn\;P_1^{-1}=sng\;P_1.$$

    {\color{def}$n$-ta grupa alternujaca} [$A_n$] - podgrupa $S_n$ zlozona ze wszystkich parzystych permutacji.
 
\end{multicols}\bigskip

\podz{sep2}\bigskip

\subsection{Grupy ilorazowe B)}
\begin{multicols}{2}
    
    {\color{def}Prawostronna warstwa} grupy $G$ wzgledem jej podgrupy $H$ wyznaczona przez $g\in G$ to zbior
    $$gH=\{gh\;:\;h\in H\},$$
    natomiast {\color{acc}lewostronna} warstwa to zbior
    $$ Hg = \{hg\;:\;h\in H\}. $$
    Dla grup abelowych sa one rowne.\medskip

    Dwa elementy $g_1,g_2\in G$ wyznaczaja te sama warstwe prawostronna wzgledem $H$, gdy $g_1^{-1}g_2\in H$, a te sama warstwe lewostronna, gdy $g_1g_2^{-1}\in H$.\medskip

    \podz{sep}\medskip

    {\color{def}Rzad grupy} skonczonej $G$ to ilosc jej elementow.\medskip

    {\color{def}Indeks} $[G : H]$ podgrupy $H$ w grupie $G$ to ilosc warstw w grupie $G$ wzgledem $H$. Dla skonczonych grup mamy:\smallskip\\
    \point $g\in G$ $o(g) | |G|$,\smallskip\\
    \point rzad i indeks kazdej podgrupy sa dzielnikami rzedu grupy,\smallskip\\
    \point jesli rzad jest liczba pierwsza, to grupa jest cykliczna\medskip

    {\color{def}Twierdzenie Lagrange'a} - dla skonczonych $G > H$:
    $$|G|=[G:H]\cdot |H|.$$

    \podz{sep}\medskip

    Podgrupa $H$ jest {\color{def}dzielnikiem normalnym} grupy $G$ [$H\triangleleft G$] jesli $(\forall\;g\in G)\;gH=Hg$. Wystarczy, ze $(\forall\;g\in G)(\forall\;h\in H)\;ghg^{-1}\in H$.\medskip

    Niech $f:G_1\to G_2$ bedzie {\color{acc}homomorfizmem}, a $e_1, e_2$ beda elementami neutralnymi grup odpowiednio $G_1,G_2$. Wtedy $f(e_1)=e_2$ oraz $f(g)^{-1}=f(g^{-1})$.\medskip

    Obraz homomorfizmu $f:G_1\to G_2$ jest {\color{acc}podgrupa grupy $G_2$} [$Im\;f < G_2$], natomiast jadro $f$ jest {\color{acc}dzielnikiem normalnym $G_1$} [$Ker\;f \dzielnat G_1$].\medskip

    \podz{sep}\medskip

    {\color{def}Grupa ilorazowa} to zbior wszyystkich warstw $H/G$, gdzie $H\dzielnat G$, z dzialaniem
    $$(g_1H)(g_2H)=(g_1g_2)H.$$
    Odwzorowanie 
    $$\phi:G\to H$$
    $$\phi(g)=gH$$
    jest {\color{acc}epimorfizmem} (czesto nazywane {\color{acc}kanonicznym homeomorfizmem $G$ na $H$}).\medskip

    {\color{cyan}[!!!]}{\color{def}Zasadnicze twierdzenie o homeomorfizmach dla grup} - jesli $f:G\to G_1$ jest epimorfizmem oraz $Ker\;f=H$, natomiast $\phi:G\to G/H$ jest dzialaniem jak wyzej, to istnieje tylko jeden izomorfizm $\psi:G/H\to G_1$ taki, ze $f=\psi\circ \phi$\medskip

    \podz{sep}\medskip

    Jezeli $\emptyset\neq A\subseteq G$ oraz $G(A)<G$ to przekroj wszystkich podgrup $G$ zawierajacych $A$, a $A\subseteq G_1<G$, to $G(A)<G_1$.\medskip

    Jezeli $K\dzielnat G$ i $H\dzielnat G$, to najmniejsza podgrupa $G$ zawierajaca $H$ i $K$ pokrywa sie ze zbiorem
    $$KH:=\{kh\;:\;k\in K,\;h\in H\}$$

    {\color{def}Pierwsze twierdzenie o izomorfizmach} - jezeli $K\dzielnat G$ i $H\dzielnat G$, to \smallskip\\
    \point $K<KH=HK<G$\smallskip\\
    \point $H\cap K\dzielnat H$ i $K\dzielnat KH$\smallskip\\
    \point $\phi:hK\to h(K\cap H)$ indukuje izomorfizm 
    $$HK/K\sim H/(H\cap K)$$

    {\color{def}Drugie twierdzenie o izomorfizmach} - jezeli $K\dzielnat G$ i $K<H<G$ i oznaczymy $\overline H = H/K$ oraz $\overline G=G/K$, to wtedy:\smallskip\\
    \point $\overline H<\overline G$\smallskip\\
    \point $\overline H\dzielnat\overline G\iff H\dzielnat G$\medskip

    \podz{sep}\medskip

    {\color{def}Automorfizm wewnetrzny }grupy $G$ wyznaczony przez $g$: $\phi_g(x)=g^{-1}xg$.\\ 
    Jesli $G$ to grupa abelowa, to dla kazdego $g$ $\phi_g(x)=x$, a wiec ma ona jedynie identycznosc.\\
    Zbior wszystkich automorfizmow wewnetrznych grupy G oznaczamy $\color{acc}I(G)$ i tworzy on grupe ze skladaniem\medskip

    {\color{def}Centrum grupy} $G$ [$Z(G)$] to zbior $x\in G$ takich, ze dla dowolnego $y\in G$ $xy=yx$. Dla kazdego $G$ $Z(G)\dzielnat G$\medskip

    Grupa $I(G)$ jest izomorficzna z $G/Z(G)$.\medskip

    Jesli $M$ to dowolny podzbior grupy $G$, to dla kazdego $g$ takiego, ze $\phi_g\in I(G)$ zbiorem {\color{def}sprzezony} do $M$ nazywamy zbior
    $$M^g=\{\phi_g(x)\;:\;x\in M\}$$
    Jesli $M=\{x\}$, to $M^g$ zawiera elementy sprzezone z x.\smallskip\\
    {\color{def}Normalizator} zbioru $M$:
    $$N_G(M)=\{g\in G\;:\;M^g=M\}$$
    {\color{def}Centralizator} zbioru $M$:
    $$C_G(M)=\{g\in G\;:\;mg=gm,\;m\in M\}$$

    Twierdzonka:\smallskip\\
    \point $(\forall\; M\subseteq G)\;C_G(M)<N_G(M)$ ($|M| = 1\implies C_G(M)=N_G(M)$)\smallskip\\
    \point $Z(G)=C_G(G)$\smallskip\\
    \point dla $M\subseteq G$ ilosc zbiorow $M^g$ jest rowna $[G\;:\;N_G(M)]$.\medskip

    Aby klasa elementow sprzezonych z $x\in G$ byla jednoelementowa wystarczy, zeby $x\in Z(G)$\medskip
    
    Jesli $G$ jest skonczona, to ilosc elementow sprzezonych z zadanym $x$ jest dzielnikiem $|G|$.\medskip

    {\color{acc}$p$-grupa} to grupa, w ktorej wszystkie elementy maja rzad $p$, gdzie $p$ jest liczba pierwsza. Jesli $|G|=p^n$ to $G$ jest $p$-grupa.\smallskip\\
    Skonczone $p$=grupy maja {\color{acc}nietrywialne centrum}.\smallskip\\
    Jesli $|G|=p^2$, to $G$ jest grupa abelowa.

\end{multicols}\bigskip

\podz{sep2}\bigskip

\subsection{Produkty grup}
\begin{multicols*}{2}
    W zbiorze $A\times B=\{(a, b)\;:\;a\in A,b\in B\}$, gdzie $A,B$ sa grupami, okreslmy
    $$(a,b)\cdot(c,d)=(ac, bd)$$
\end{multicols*}

\newpage

\section{External problems}

\pdef

    \emph{How large can we make some parameter of $G$ before it is forced to have a certain property?}

\kdef

\subsection{Complete subgraphs}

\begin{tabularx}{\textwidth}{ X!{\color{git90gray}\vrule} X }

{\color{def}Complete graph of order $r$} [$K_r$] is a graph with $V(K_r)=[r]$ and $E(K_r)=\{ij\;:\;1\leq i<j\leq r\}$.

$K_3$ is called a {\color{acc}triangle}.
\medskip

{\color{def}$r$-partite} graph $G$ with vertex classes $V_1,...,V_r$ has every edge of form $vw$ where $v\in V_i$ and $w\in V_j$ and $i\neq j$. Such a graph is called {\color{acc}complete $r$-partite} if for every $i,j,\;i\neq j$ we have $v\in V_i,\;w\in V_j\implies vw\in E(G)$.
\medskip

A {\color{def}complete bipartite graph} with vertex classes of orders $|V_1|=m$ and $|V_2|=n$ is denoted $K_{m,n}$

&

{\color{def}Graf pelny stopnia $r$} [$K_r$] to graf z $V(K_r)=[r]$ i $E(K_r)=\{ij\;:\;1\leq i<j\leq r\}$.

$K_3$ jest nazywany {\color{dyg}trojkatem}?
\medskip

{\color{def}$r$-dzielny} graf $G$ z klasami wierzcholkow $V_1,...,V_r$ ma kazdy wierzcholek postaci $vw$, gdzie $v\in V_i$ i $w\in V_j$ dla $i\neq j$. Taki graf jest dodatkowo nazywany {\color{acc}pelnym grafem $r$-dzielnym}, jezeli dla kazdego $i,j,\;i\neq j$ jest $v\in V_i,\;w\in V_j\implies vw\in E(G)$.
\medskip

{\color{def}Pelny graf dwudzielny} z klasami wierzchoklow o mocy $|V_1|=m$ i $|V_2|=n$ jest oznaczany jako $K_{m,n}$.
\\

& \\

\hline

& \\

We now want to check how big must \hyperref[handshaking-lemma]{$e(G)$} be in order to force $K_r$ to be $G$ subgraph.
\medskip

{\color{acc}\point }
given $r\geq 2$ we can see that for $G$ to be $K_r$-free we need $G$ to be $(r-1)$-partite
\smallskip

{\color{acc}\point }
given $n\geq r-1$ out of all $(r-1)$-partite graphs with $n$ vertices the one with most edges is a complete $(r-1)$-partite graph
\smallskip

{\color{acc}\point }
if $G$ is a complete $(r-1)$-partite graph with vertex classes 

$V_1,..., V_{r-4}$, 

if $|V_i|\geq |V_j|+2$ for some $i\neq j$, then we may choose $v\in V_j$ and consider a new graph obtained by removing edges $vv_i$ for $v_i\in V_i$ and adding $vv_j$ for every $v_j\in V_j-\{v\}$. This new graph is $(r-1)$-partite and $|G'|=|G|$ and 
$$e(G')=e(G)-|V_i|+|V_j|-1>e(G)$$

{\color{acc}\point} so the $(r-1)$ partite graph with $n$ vertices and the most edges will have vertes classes "as equal in size as possible"

&

Teraz chcemy sprawdzic jak duze musi byc \hyperref[handshaking-lemma]{$e(G)$}, zeby zmusic $K_r$ do bycia podgrafem $G$
\medskip

{\color{acc}\point} majac dane $r\geq2$ mozemy zauwazyc, ze aby $G$ bylo $K_r$-wolne, musi byc $(r-1)$-dzielne
\smallskip

{\color{acc}\point} majac dane $n\geq r-1$ i wszystkie $(r-1)$-dzielne grafy z $n$ wierzcholkammi ten, ktory ma najwiecej krawedzi jest pelnym $(r-1)$-dzielnym grafem
\smallskip

{\color{acc}\point} jesli $G$ jest grafem pelnym $(r-1)$-dzielnym z klasami wierzcholkow 

$V_1,...,V_{r-1}$,

to jesli $|V_i|\geq |V_j|+2$ dla pewnych $i\neq j$, wtedy mozemy wybrac $v\in V_j$ i rozwazyc nowy graf otrzymany poprzez usuniecie krawedzi $vv_i$ dla $v_i\in V_i$ oraz dodanie $vv_j$ dla kazdego $v_j\in V_j-\{v\}$. Taki nowy graf jest $(r-1)$-dzielny i $|G'|=|G|$ oraz
$$e(G')=e(G)-|V_i|+|V_j|-1>e(G)$$

{\color{acc}\point}wiec $(r-1)$-dzielny graf z $n$ wierzchoklami i najwieksza liczba krawedzi bedzie mial klasy wierzcholkow tak bliskie rozmiarem jak mozliwe
\\

& \\

\hline

& \\

{\color{def}Tur\'an graph} $T_r(n)$ for $n\geq r\geq 1$ is a complete $r$-partite graph of order $n$ with all vertex classes of size $\lfloor{n\over r}\rfloor$ or $\lceil{n\over r}\rceil$. We denote $\color{acc}e(T_r(n))$ as $t_r(n)$.

&

{\color{def}Graf Tur\'ana} $T_r(n)$ dla $n\geq r\geq 1$ to pelny $r$-dzielny graf stopnia $n$ ze wszystkimi klasami wierzcholkow roozmiaru $\lfloor{n\over r}\rfloor$ lub $\lceil{n\over r}\rceil$. Oznaczamy $\color{acc}e(T_r(n))$ jako $t_r(n)$.

\end{tabularx}
\bigskip

\begin{tabularx}{\textwidth}{ X!{\color{git90gray}\vrule} X }

Some observations:
\medskip

{\color{acc}\point} If $T_r(n)\simeq G-\{e\}$ for some $e\in E(G)$, then $G$ is {\color{acc}not $K_{r+1}$-free}.
\medskip

{\color{acc}\point} If {\color{acc}$r$ divides $n$}, then we have
$$\delta(T_r(n))=d(T_r(n))=\Delta(T_r(n))=n-\frac nr,$$

otherwise vertices in large classes have {\color{acc}minimal degree}, $\delta(T_r(n))=n-\lceil\frac nr\rceil$, and vertices in small classes have {\color{acc}maximal degree}, $\Delta(T_r(n))=n-\lfloor\frac nr\rfloor$. This implies that always
$$\delta(T_r(n))=\lfloor d(T_r(n))\rfloor$$
$$\Delta(T_r(n))=\lceil d(T_r(n))\rceil$$

{\color{acc}\point} $T_r(n-1)\simeq T_r(n)-\{v\}$ where $v\in T_r(n)$ is a vertex of minimal degree (any if $r|n$ or one of the vertices in large classes).
\medskip

{\color{acc}\point} $t_r(n-1)=t_r(n)-\delta(T_r(n))$
\medskip

{\color{acc}\point} Let us say that we want to create a {\color{def}$K_{r+1}$-free graph $G$ by adding a vertex $v$} to $T_r(n-1)$ and $m$ edges with $m$ being as large as posiible. We know that $v$ cannot be adjavent to a vertex in every class so we have
$$m=d_G(v)\leq n-1-\lfloor{n-1\over r}\rfloor=n-\lceil{n\over r}\rceil,$$
with {\color{acc}equality} iff $G$ is complete $r$-partite (obtained by adding $v$ anywhere if $r|(n-1)$ or to one of the small classes.)

&

Kilka obserwacji:
\medskip

\phantomsection
\label{turan-observ1-PL}
{\color{acc}\point} Jesli $T_r(n)\simeq G-\{e\}$ dla pewnego $e]in E(G)$, wtedy $G$ {\color{acc}nie jest $K_{r+1}$-free}{\color{dyg}????}.
\medskip

\phantomsection
\label{turan-observ-PL}
{\color{acc}\point} Jesli {\color{acc}$r$ dzieli $n$}, to wtedy
$$\delta(T_r(n))=d(T_r(n))=\Delta(T_r(n))=n-\frac nr,$$

w przeciwnym wypadku wierzcholki w duzych klasach maja {\color{acc}najmniejszy stopien}, $\delta(T_r(n))=n-\lceil\frac nr\rceil$, i wierzcholki w malych klasach maja {\color{acc}najwiekszy stopien}, $\Delta(T_r(n))=n-\lfloor\frac nr\rfloor$. To implikuje, ze zawsze
$$\delta(T_r(n))=\lfloor d(T_r(n))\rfloor$$
$$\Delta(T_r(n))=\lceil d(T_r(n))\rceil$$

\phantomsection
\label{turan-observ2-PL}
{\color{acc}\point} $T_r(n-1)\simeq T_r(n)-\{v\}$ gdzie $v\in T_r(n)$ jest wierzcholkiem najmniejszego stopnia (kazdy jesli $r|n$, wpp jeden z wierzcholkow w duzych klasach).
\medskip

{\color{acc}\point} $t_r(n-1)=t_r(n)-\delta(T_r(n))$
\medskip

\phantomsection
\label{turan-observ3-PL}
{\color{acc}\point} Powiedzmy, ze chcemy stworzyc {\color{def}$K_{r+1}$-free graf $G$ poprzez dodanie wierzcholka $v$} do $T_r(n-1)$ i $m$ krawedzi, gdzie $m$ jest najwieksze mozliwe. Wiemy, ze $v$ nie moze byc obok wierzcholkow w kazdej klasie, wiec
$$m=d_G(v)\leq n-1-\lfloor{n-1\over r}\rfloor=n-\lceil{n\over r}\rceil,$$
z {\color{acc}rownoscia} iff $G$ jest pelnym grafem $r$-dzielnym (otrzymanym przez dodanie $v$ gdziekolwiek jesli $r|(n-1)$ lub do jednej z malych klas).
\\

& \\

\hline

& \\

{\color{def}Tur\'an's Theorem}

Let $n\geq r\geq 1$ and let $G$ be a $K_{r+1}$-free graph with $|G|=n$ and $e(G)\geq t_r(n)$. Then $\color{acc}G\simeq T_r(n)$.

&

{\color{def}Twierdzenie Tur\'ana}

Niech $n\geq r\geq 1$ i niech $G$ bedzie $K_{r+1}$-wolnym grafem z $|G|=n$ i $e(G)\geq t_r(n)$. Wtedy $\color{acc} G\simeq T_r(n)$.

\end{tabularx}

\prooflink{turan-theo}
\medskip

\prooflang{turan-theo}{GB}
\medskip

{\color{cyan}Maybe later, dunno}
\medskip

\prooflang{turan-theo}{PL}
\medskip

Uzyjemy indukcji na $n$.
\smallskip

Jesli $n=r$, to $T_r(n)\simeq K_r$ i mamy
$${n\choose 2}=t_r(n)\leq e(G)\leq {n\choose 2}$$
i z tego $T_r(n)\simeq G$.
\smallskip

Niech teraz $n>r$. Wybierzmy podzbior $E'\subseteq E(G)$ taki, ze $|E'|=e(G)-t_r(n)$ i niech $H=G-E'$, czyli $e(H)=t_r(n)$. Wtedy mamy
$$d(H)={2e(H)\over n}={2t_r(n)\over n}=d(T_r(n)).$$
Wtedy 
$$\delta(H)\leq \lfloor d(H)\rfloor=\lfloor d(T_r(n))\rfloor=\delta(T_r(n))$$
gdzie ostatnia rownosc wynika z \hyperref[turan-observ-PL]{\color{dyg}obserwacji}.

Wybierzmy teraz $v\in H$ taki, ze $d(v)=\delta(H)$ i niech 
$$K=H-\{v\}.$$
Wtedy $K$ jest $K_{r+1}$-wolne. $|K|=n-1$ i
$$e(K)=e(H)-d_H(v)=t_r(n)-\delta(H)\geq t_r(n)-\delta(T_r(n))=t_r(n-1),$$
gdzie ostatnia rownosc wynika z \hyperref[turan-observ2-PL]{\color{dyg}obserwacji}. W takim razie, poprzez hipoteze indukcyjna, wiemy, ze
$$K\simeq T_r(n-1).$$
W szczegolnosci, z tego wynika, ze
$$e(K)=t_r(n-1)$$
i w takim razie
$$d_H(v)=\delta(T_r(n))$$
wiec z kolejnej \hyperref[turan-observ3-PL]{\color{dyg}obserwacji} mamy, ze $H\simeq T_r(n)$.

W koncu, poniewaz $V(H)=V(G)$ oraz $E(H)=E(G)-E'\subseteq E(G)$ i $G$ jest $K_{r+1}$-wolne, to z \hyperref[turan-observ1-PL]{\color{dyg}obserwacji} mamy $|E'|=0$ i
$$G\simeq H\simeq T_r(n).$$

\subsection{Complete bipartite subgraphs}

\begin{tabularx}{\textwidth}{ X!{\color{git90gray}\vrule} X }

    {\color{def}Jensen's Inequality}: $a<b\in\R$ and $f:[a,b]\to\R$ is convex. Then
    $$\frac1n\sum\limits_{i=1}^nf(x_1)\geq f(\frac1n\sum\limits_{i=1}^nx_i)$$
    for all $x_1,...,x_n\in[a,b]$.

\end{tabularx}

\subsection{Arbitrary subgraphs}

\begin{tabularx}{\textwidth}{ X!{\color{git90gray}\vrule} X }

    {\color{def}Forbidden subgraph problem}: given a graph $H$, how many edges can a $H$-free graph of order $n$ have?
    $$\color{def}ex(n; H)=\max\{e(G)\;:\;\text{G - H-free graph with }|G|=n\}.$$

    From previous theorems we know that:\smallskip

    \point $ex(n; K_{r+1})=t_r(n)$ and $ex(n;K_{r+1})\sim {n^2\over2}(1-\frac1r)$

    \point $ex(n; K_{t,t})=O(n^{2-\frac1t})$
    \medskip

    Let us write 
    $$\color{acc}e(H)=\lim_{n\to\infty}{ex(n;H)\over{n\choose 2}}$$

    and proof that for $n\geq 2$ the sequence
    $$({ex(n;H)\over{n\choose 2}})_{n=2}^\infty$$
    converges.

    &

    {\color{def}Problem zakazanego podgrafu}: mając dany graph $H$, ile krawędzi może mieć $H$-wolny graf rzędu $n$?
    $$\color{def}ex(n; H)=\max\{e(G)\;:\;\text{G - H-wolny graf z }|G|=n\}.$$

    Z poprzednich twierdzeń wiemy, że:\smallskip

    \point $ex(n; K_{r+1})=t_r(n)$ oraz $ex(n;K_{r+1})\sim {n^2\over2}(1-\frac1r)$

    \point $ex(n; K_{t,t})=O(n^{2-\frac1t})$
    \medskip

    Oznaczmy
    $$\color{acc}e(H)=\lim_{n\to\infty}{ex(n;H)\over{n\choose 2}}$$

    i udowodnijmy, że dla $n\geq 4$, ciąg
    $$({ex(n;H)\over{n\choose 2}})_{n=2}^\infty$$
    jest zbieżny.    

\end{tabularx}

\prooflink{forbidden-convergence}
\medskip

\prooflang{forbidden-convergence}{GB}
\medskip

DUNNO
\medskip

\prooflang{forbidden-convergence}{PL}
\medskip

Ciąg $(x_n)_{n=2}^\infty$ jest ograniczony od dołu przez $0$, więc wystarczy pokazać, że nie jest to ciąg rosnący. Niech $n\geq3$ i niech $G$ będzie $H$-wolnym grafem z $|G|=n$ oraz $e(G)=ex(n; H)$. Wtedy dla dowolnego $v\in G$ graf $G-\{v\}$ jest $H$-wolny i ma rząd $n-1$, implikując, że
$$e(G-\{v\})\leq ex(n-1, H).$$

Z drugiej strony, dowolna krawędź $uw\in E(G)$ jest używa t dokładnie $n-2$ grafach $G-\{v\}$ dla $v\in G$ (w tych, gdzie $v\notin\{u,w\}$). To z kolei implikuje, że
$$(n-2)e(G)=\sum\limits_{v\in G}e(G-\{v\})$$
i z tego mamy
$$x_n={ex(n;H)\over {n\choose2}}={2e(G)\over n(n-1)}=\sum\limits_{v\in G}{2e(G-\{v\})\over n(n-1)(n-2)}\leq {2ex(n-1;H)\over (n-1)(n-2)}$$
czyli ciąg nie jest rosnący, a więc musi do czegoś zbiegać.

\end{document}